\documentclass[10pt,a4paper]{article}
\usepackage[T1]{fontenc}
\usepackage[utf8]{inputenc}
\usepackage[margin=2.2cm]{geometry}
\usepackage[strings]{underscore}
\usepackage{xcolor}
\usepackage{hyperref}
\usepackage{titlesec}
\usepackage{microtype}

\definecolor{swgrey}{gray}{0.6}
\setlength{\parindent}{0pt}
\titleformat{\section}{\Large\bfseries}{}{0pt}{}
\titleformat{\subsection}{\large\bfseries\color{black!70}}{}{0pt}{}

% \sw{name}{type}{default}{description}
\newcommand{\sw}[4]{%
  \par\vspace{0.7em}%
  \begin{minipage}{\linewidth}\raggedright
    \textbf{\texttt{#1}}\hfill\textit{#2}\\[0.1em]
    \textbf{Default:} #3\\[0.1em]
    #4\par
    \vspace{0.3em}{\color{swgrey}\hrule}
  \end{minipage}\par}

\title{GOAT Input Switches Reference}
\author{Generated from \texttt{goatmod\_userinput}}
\date{}

\begin{document}
\maketitle
\tableofcontents

\section*{Notes}
\begin{itemize}
  \item Each switch is listed under the key as it appears in the options file (e.g.\ \texttt{goat.debug}). Keys are read by the read routines in the \texttt{goatmod\_userinput} module.
  \item Types: \emph{integer (0 or 1)} (\texttt{1}/\texttt{0}), \emph{integer}, \emph{real}, \emph{string}, and the array variants (\emph{integer array}, \emph{real array}). Numbers provided in the arrays \textbf{must} be separated by a comma.
  \item "Empty" means a zero-length array.
  \item When SOLPS mode is active, the I/O file paths default to the \texttt{solps\_*} paths of \texttt{mod\_global\_environment}. This is in principle available but it is not yet fully supported.
  \item Discaimer: the switches \texttt{gg.ref.XXX} are meant for poloidal refinement, while \texttt{gg.radref.XXX} are meant for radial refinement. Ideally speaking, \texttt{gg.ref.XXX} would be equilvalent to \texttt{gg.polref.XXX}.
\end{itemize}

%======================================================================
\section{GOAT -- General}
%======================================================================
\subsection*{General}
\sw{goat.debug}{integer (0 or 1)}{\texttt{0}}{If set to \texttt{1} and the \texttt{goat_debug} execuatable is run, it provides extra screen output.}
\sw{goat.meth}{string}{\texttt{`GG'}}{GOAT running method. Currently, only \texttt{`GG'} (Grid Generation) is supported. Users can select \texttt{`GD'} for Grid Deformation or \texttt{`GA'} for Grid Adaptation, yet those are not fully mature/supported.}

%======================================================================
\section{GOAT -- Environment}
%======================================================================
\sw{goat.readstate}{integer (0 or 1)}{\texttt{0}}{Read the plasma state.}
\sw{goat.readstatemeth}{string}{\texttt{`b2fplasmf'}}{Method to read the state; can be \texttt{`b2fstate'} or \texttt{`b2fplasmf'}.}
\sw{goat.statefilepath}{string}{\texttt{`./b2fplasmf'} (SOLPS: \texttt{solps\_statefilepath})}{State file to read.}


\subsection*{Input file names}
\sw{goat.gridfilepath}{string}{\texttt{`./traduit.out.b2us'} (SOLPS: \texttt{solps\_gridfilepath})}{File path to the file with grid data (e.g.\ the \texttt{traduit.out.b2us} file).}
\sw{goat.structurefilepath}{string}{\texttt{`./structure.dat'} (SOLPS: \texttt{solps\_structurefilepath})}{The \texttt{structure.dat} file to read.}
\sw{goat.magneticfieldfilepath}{string}{\texttt{`./rzpsi.dat'} (SOLPS: \texttt{solps\_magneticfieldfilepath})}{Magnetic field file to read.}
\sw{goat.GDinputfilepath}{string}{\texttt{`./GOAToptions.dat'}}{Path towards the file where the options for grid deformation are defined.}

\subsection*{Output options}
\sw{goat.writefilepath}{string}{\texttt{`traduit.out.b2us'} (SOLPS: \texttt{solps\_writefilepath})}{Path where the output traduit file is written.}
\sw{goat.write_final}{integer (0 or 1)}{\texttt{1}}{Write final output.}
\sw{goat.write_traduitb2us}{integer (0 or 1)}{\texttt{1}}{Write unstructured traduit file.}
\sw{goat.write_b2agdat}{integer (0 or 1)}{\texttt{1}}{Write final \texttt{b2ag.dat} file for use in b2ag.}
\sw{goat.write_OMPdata}{integer (0 or 1)}{\texttt{0}}{Write OMP data in the traduit file.}
\sw{goat.write_topologicaldata}{integer (0 or 1)}{\texttt{0}}{Write out X-, O-, and strike point information, and other topological data.}

\subsection*{Case identification (TO BE CHECKED)}
\sw{goat.vesselmode}{integer (0 or 1)}{\texttt{0}}{Set to true if the case is a vessel mode grid.}
\sw{goat.slab}{integer (0 or 1)}{\texttt{0}}{True if slab grid.}
\sw{goat.artificial_slab}{integer (0 or 1)}{\texttt{0}}{True if artificial slab.}

\subsection*{Face label mappings}
Those are only required when working with GD and GA.
\sw{goat.facelabelmappingGG}{integer array}{empty}{Labels as defined in the grid generator. Must have the same length as \texttt{goat.facelabelmappingGA}.}
\sw{goat.facelabelmappingGD}{integer array}{empty}{Corresponding labels for GD (the first GG label is mapped to the first GD label here).}
\sw{goat.facelabelmappingGA}{integer array}{empty}{Corresponding labels for GA (the first GG label is mapped to the first GA label here).}
\sw{goat.facelabelsubfrom}{integer array}{empty}{Substitution of this face label \dots}
\sw{goat.facelabelsubto}{integer array}{empty}{\dots\ to this face label in the GA to GD interface.}

\subsection*{OMP and IMP definition}
\sw{goat.OMP_r}{real array (2)}{\texttt{[0, 0]}}{R coordinates that define the outer mid plane line segment.}
\sw{goat.OMP_z}{real array (2)}{\texttt{[0, 0]}}{Z coordinates that define the outer mid plane line segment.}
\sw{goat.IMP_r}{real array (2)}{\texttt{[0, 0]}}{R coordinates that define the inner mid plane line segment.}
\sw{goat.IMP_z}{real array (2)}{\texttt{[0, 0]}}{Z coordinates that define the inner mid plane line segment.}


%======================================================================
\section{GOAT -- Grid}
%======================================================================
\sw{goat.grid.type}{string}{\texttt{`plasma'}}{The type of grid that is considered. Only \texttt{`plasma'} is currently available.}
\sw{goat.grid.readmeth}{string}{\texttt{`b2fgmtry'}}{The type of input file. Can be \texttt{`b2fgmtry'} or \texttt{`traduitb2us'}. Both are unstructured.}

%======================================================================
\section{GOAT -- Magnetic Field}
%======================================================================
\sw{goat.mf.readmeth}{string}{\texttt{`default'}}{Method to read in the magnetic field data. Only \texttt{`readrzpsi'} and \texttt{`default'} are supported for now; both use the \texttt{readrzpsi} routine.}
\sw{goat.mf.interpC}{integer}{\texttt{3}}{Desired continuity of the interpolant.}
\sw{goat.mf.interpM}{integer}{\texttt{6}}{Order of the interpolant describing the values at the grid nodes.}
\sw{goat.mf.interpmeth}{string}{\texttt{`uniformgrid'}}{Interpolant method: \texttt{`uniformgrid'} for data given on a uniform grid, \texttt{`centered'} for default non-uniform grids.}
\sw{goat.mf.reinterpolate}{integer (0 or 1)}{\texttt{0}}{Reinterpolate the magnetic field for a different resolution defined by \texttt{goat.mf.resx} and \texttt{goat.mf.resy}.}
\sw{goat.mf.resx}{integer}{\texttt{100}}{Resolution in x-direction for reinterpolation.}
\sw{goat.mf.resy}{integer}{\texttt{100}}{Resolution in y-direction for reinterpolation.}
\sw{goat.mf.RBtor}{real}{\texttt{0}}{Product of toroidal magnetic field and major radius (assumed constant here) in [Tm]. Note: this value can be signed, the sign indicating positive or negative toroidal direction of the magnetic field.
The direction is not important for GOAT purposes, and it is completely related to the final application of the grid: e.g., for a SOLPS-ITER grid, a positive value provided here will end up in a positive B field in SOLPS-ITER coordinate system.}

%======================================================================
\section{GOAT -- Vessel}
%======================================================================
\sw{goat.vessel.readmeth}{string}{\texttt{`read\_structure'}}{\texttt{'read\_structure'} for \texttt{structure.dat} files. This is the only method currently supported.}
\sw{goat.vessel.exclude}{integer array}{empty}{List of structures in the vessel to be excluded from the vessel (indices must be the same as in \texttt{structure.dat}; those are also excluded when generating the bounding polygon).}
\sw{goat.vessel.refinevessel}{integer (0 or 1)}{\texttt{0}}{Set to true to refine the vessel (insert more nodes).}
\sw{goat.vessel.maxvesseldist}{real}{\texttt{0.01}}{Maximum distance between two nodes. If larger, nodes are added in between when refinement is enabled.}
\sw{goat.vessel.minreffac}{integer}{\texttt{0}}{Minimal refinement factor for each edge. If zero, no effect. The value plus one gives the minimal number of edges into which each edge is split.}
\sw{goat.vessel.TP}{integer array}{\texttt{[1, 2]}}{Indices of which part of the structure should be considered as target plates. This should not account for structures excluded with \texttt{goat.vessel.exclude}.}
\sw{goat.vessel.TPind}{integer array}{\texttt{[1, 2]}}{Target plate enumerators.}
\sw{goat.vessel.shapemeth}{string}{\texttt{`closedpolygon\_smoothapproximation'}}{Method used to represent (and possibly approximate) the shape of the vessel.}
\sw{goat.vessel.resx}{integer}{\texttt{100}}{Resolution in x-direction for the shape representation.}
\sw{goat.vessel.resy}{integer}{\texttt{100}}{Resolution in y-direction for the shape representation.}
\sw{goat.vessel.interpC}{integer}{\texttt{3}}{Only used for interpolant-based methods. Continuity order of the interpolant.}
\sw{goat.vessel.interpM}{integer}{\texttt{6}}{Only used for interpolant-based methods. Order used to approximate derivatives to construct the interpolant function.}
\sw{goat.vessel.offsetfracx}{real}{\texttt{0.1}}{Fractional offset to be taken from the minimal and maximal x-value of the original vessel polygon.}
\sw{goat.vessel.offsetfracy}{real}{\texttt{0.1}}{Same as \texttt{offsetfracx}, but for the y-direction.}
\sw{goat.vessel.xrange}{real array}{empty}{Only used for interpolant-based methods. Range over which the interpolant can be evaluated (determined by max and min value in the array; may be overwritten).}
\sw{goat.vessel.yrange}{real array}{empty}{Same as \texttt{xrange}, for the y-direction.}

\newpage
%======================================================================
\section{Grid Generation (GG)}
%======================================================================
%======================================================================
\subsection*{General}
%======================================================================
\sw{gg.verbosity}{integer}{\texttt{1}}{The higher, the more information is printed (0 suppresses all).}
\sw{gg.vertexconstructionmethod}{string}{\texttt{`independent'}}{How to construct the grid: \texttt{`independent'} treats each flux surface independently; \texttt{`orthogonal'} yields a largely orthogonal grid but needs dependency between mesh cells.}
\sw{gg.cellconstructionmethod}{string}{\texttt{`quads\_triangles'}}{Method to determine grid cells starting from a given vertex distribution. \texttt{`quads\_triangles'} is recommended.}
\sw{gg.legalcellstyle}{string}{\texttt{`graph'}}{How to check that cells are legal. \texttt{`no'}: no checks; \texttt{`old'}: deprecated method that catches most issues but not all; anything else: graph-based method, likely more expensive but captures all cell-overlap cases.}
\sw{gg.TMcellgriddingorder}{string}{\texttt{`sequential'}}{Order in which to grid the topological mesh cells: \texttt{`independent'} (essentially random) or \texttt{`sequential'} (recommended; initial distributions are propagated through).}
\sw{gg.dogriddiagnostics}{integer (0 or 1)}{\texttt{1}}{Run grid diagnostics (time consuming).}

\subsection*{Topological Mesh (Topomesh)}
\subsubsection*{I/O}
\sw{gg.tm.readexistingTM}{integer (0 or 1)}{\texttt{0}}{Read an existing topomesh file, whose full path is defined in \texttt{gg.tm.TMfilepath}.}
\sw{gg.tm.readexistingtracers}{integer (0 or 1)}{\texttt{0}}{Read existing field and vessel tracers instead of constructing new ones; only when restarting the topomesh from a previous one. Paths are given in \texttt{gg.tm.TMfieldtracerfilepath} and \texttt{gg.tm.TMvesseltracerfilepath}.}
\sw{gg.tm.writedebugoutput}{integer (0 or 1)}{\texttt{0}}{Write (many) more intermediate topomesh files for visual debugging.}
\sw{gg.tm.TMfilepath}{string}{\texttt{`topomesh.dat'}}{Full path of the topomesh file to read.}
\sw{gg.tm.TMfieldtracerfilepath}{string}{\texttt{`./output/TMfieldtracer.dat'}}{Full path of the field tracer file.}
\sw{gg.tm.TMvesseltracerfilepath}{string}{\texttt{`./output/TMvesseltracer.dat'}}{Full path of the vessel tracer file.}

\subsection*{Resolution and contouring}
\sw{gg.tm.field.resx}{integer}{\texttt{100}}{Resolution in x for 2D tracing of contours and extrema of the field.}
\sw{gg.tm.field.resy}{integer}{\texttt{100}}{Resolution in y for 2D tracing of contours and extrema of the field.}
\sw{gg.tm.vessel.resx}{integer}{\texttt{100}}{Same as field resolution, but for the vessel (x).}
\sw{gg.tm.vessel.resy}{integer}{\texttt{100}}{Same as field resolution, but for the vessel (y).}
\sw{gg.tm.contour.npmin}{integer}{\texttt{10}}{Minimal number of points of contours (when coarsening).}
\sw{gg.tm.contour.npmax}{integer}{\texttt{10000}}{Maximal number of points of contours (when coarsening).}
\sw{gg.tm.contour.dl}{real}{\texttt{1e-3}}{Desired uniform edge length.}

\subsection*{Refinement of extrema}
\sw{gg.tm.field.donewton}{integer (0 or 1)}{\texttt{1}}{Refine extrema of the field with a Newton solver (may not always converge).}
\sw{gg.tm.vessel.donewton}{integer (0 or 1)}{\texttt{1}}{Same, for the vessel.}
\sw{gg.tm.fieldtol}{real}{\texttt{1e-8}}{Tolerance on field value of extrema: if the difference is below the tolerance, two extrema are considered to have exactly the same field value.}
\sw{gg.tm.dotpvesselbased}{integer (0 or 1)}{\texttt{0}}{Do tangency point determination purely based on the current vessel polygons.}

\subsection*{Adaptations}
\sw{gg.tm.doadaptations}{integer (0 or 1)}{\texttt{1}}{General switch to apply adaptations to the basic topological mesh (core/PF boundaries, optional region removal, tube merging, \dots).}
\sw{gg.tm.addcoreboundaries}{integer (0 or 1)}{\texttt{1}}{Add additional core boundaries.}
\sw{gg.tm.addPFboundaries}{integer (0 or 1)}{\texttt{1}}{Add additional private-flux-like boundaries (at tangency points of which the contour does not enter the vessel).}
\sw{gg.tm.removecoreregions}{integer (0 or 1)}{\texttt{1}}{Remove the innermost core regions.}
\sw{gg.tm.removewidegridregions}{integer (0 or 1)}{\texttt{1}}{Remove all regions that are not next to a separatrix.}
\sw{gg.tm.removenoncoreregions}{integer (0 or 1)}{\texttt{0}}{Remove non-core regions.}
\sw{gg.tm.cbnd.frac}{real}{\texttt{0.2}}{Fraction in field value between the core and the connecting X-point (1: tangency point, 0: other PF boundary, e.g.\ separatrix).}
\sw{gg.tm.PFbnd.frac}{real}{\texttt{0.2}}{Fraction in field value between the outer PF boundary and the tangency point.}

\subsection*{Tube merging and separatrix simplification}
\sw{gg.tm.mergetubes}{integer (0 or 1)}{\texttt{0}}{Apply general tube merging.}
\sw{gg.tm.mergetubes.meth}{string}{\texttt{`extensive'}}{Method to merge tubes: \texttt{`extensive'} considers several tubes at once and often gives nicer tubes, but may fail to merge when cycles are present in the tube graph; \texttt{`local'} considers only two tubes at a time and can merge more pairs, but may lead to `uglier' tubes.}
\sw{gg.tm.mergetubes.allowseparatrix}{integer (0 or 1)}{\texttt{0}}{Allow merging over the separatrix in the general merge case.}
\sw{gg.tm.mergetubes.allowcore}{integer (0 or 1)}{\texttt{0}}{Same as for the separatrix, but for core contours.}
\sw{gg.tm.mergetubes.allowpf}{integer (0 or 1)}{\texttt{0}}{Same as for the separatrix, but for PF contours.}
\sw{gg.tm.mergetubes.oldstyle}{integer (0 or 1)}{\texttt{0}}{Use old-style tube merging.}
\sw{gg.tm.dpsimintubes}{real}{\texttt{0.0}}{Minimal psi difference for topomesh tubes, below which merging is attempted. Zero to ignore.}
\sw{gg.tm.lradmintubes}{real}{\texttt{0.0}}{Minimal radial length for topomesh tubes, below which merging is attempted. Zero to ignore. (Read twice in the source; harmless.)}
\sw{gg.tm.simplifyseparatrix}{integer (0 or 1)}{\texttt{0}}{Simplify the separatrix (remove short separatrix faces).}
\sw{gg.tm.simplifyseparatrix.lpolsepmin}{real}{\texttt{0.0}}{Minimal separatrix face length in poloidal direction. If lower, separatrix faces can be removed during the separatrix simplification step.}

\subsection*{Vessel region removal}
\sw{gg.tm.removevesselregions}{integer (0 or 1)}{\texttt{0}}{Remove regions that are adjacent to some user-specified vessel region.}
\sw{gg.tm.rvrvesselIDs}{integer array}{empty}{IDs, as specified in the \texttt{structure.dat} or \texttt{vessel.dat} file, to consider.}
\sw{gg.tm.rvrfullycovered}{integer (0 or 1)}{\texttt{0}}{If true, a face is only considered if it is fully covered by one (or multiple) defined vessel IDs.}
\sw{gg.tm.rvrretain}{integer (0 or 1)}{\texttt{0}}{Retain tubes instead of deleting them, and delete all others.}
\sw{gg.tm.rvrdocascade}{integer (0 or 1)}{\texttt{0}}{Delete not only marked tubes, but also tubes that connect to them (cascade the deletion) in the direction set by \texttt{rvrcascadedir}.}
\sw{gg.tm.rvrcascadedir}{string}{\texttt{`none'}}{\texttt{`upwards'} for tubes with higher psi value, \texttt{`downwards'} for lower, \texttt{`none'} for no cascade.}

\subsection*{Aligned vessel parts (AVP)}
\sw{gg.tm.alignvesselparts}{integer (0 or 1)}{\texttt{0}}{Define certain vessel parts as aligned faces with a certain flux surface value and ID. Only certain boundary faces are considered (typically those near type 2 tangency points). Done only after the initial topological mesh construction, as it triggers profound adaptation of the topomesh.}
\sw{gg.tm.avpminangle}{real}{\texttt{0.0}}{Minimum angle (in degrees) of boundary edges w.r.t.\ the magnetic field. If below, the edge is considered a potential aligned part.}
\sw{gg.tm.avprefinevessel}{integer (0 or 1)}{\texttt{0}}{(Deprecated) Refine vessel boundaries, similar to full vessel refinement.}
\sw{gg.tm.avpmaxvesseldist}{real}{\texttt{1e-2} [m]}{(Deprecated) Maximal vessel edge length.}
\sw{gg.tm.avpminreffac}{integer}{\texttt{0}}{(Deprecated) Minimal refinement factor for vessel edge refinement.}

\subsubsection*{Label translation}
\sw{gg.labels.structurebased}{integer (0 or 1)}{\texttt{0}}{Base labels on structure IDs.}
\sw{gg.labels.forceSOLPStopology}{integer (0 or 1)}{\texttt{0}}{Force a certain topology for region translation etc.}
\sw{gg.labels.SOLPStopology}{string}{empty string}{Desired topology: \texttt{`general'}, \texttt{`linear'}, \texttt{`limiter'}, \texttt{`single\_null'}, \texttt{`double\_null'}.}

\subsubsection*{Contouring and streamline tracing}
\sw{gg.vd.contouring.resx}{integer}{\texttt{100}}{Contouring resolution in x within the grid generator.}
\sw{gg.vd.contouring.resy}{integer}{\texttt{100}}{Contouring resolution in y within the grid generator.}
\sw{gg.vd.contouring.coarsen}{integer (0 or 1)}{\texttt{0}}{Coarsen contours.}
\sw{gg.vd.orthlinetracing.step}{real}{\texttt{0.5}}{Step size of the orthogonal streamline tracer.}
\sw{gg.vd.orthlinetracing.nsteps}{integer}{\texttt{2000}}{Number of steps of the orthogonal streamline tracer.}

\subsubsection*{Poloidal vertex distribution}
\sw{gg.vd.pd.type}{string}{\texttt{`densitybased'}}{\texttt{`uniform'} for a uniform distribution; \texttt{`densitybased'} for a non-uniform distribution based on a predefined node density distribution function.}
\sw{gg.vd.pd.distribution.plftype}{string}{\texttt{`target'}}{Type of attraction at the plasma-facing location for the density-based distribution.}
\sw{gg.vd.pd.distribution.facelength}{real}{\texttt{4e-2}}{Face length for the \texttt{`uniform'} distribution.}
\sw{gg.vd.pd.distribution.decaylengthplf}{real}{\texttt{0.05}}{Decay length from the vessel for the \texttt{`densitybased'} option.}
\sw{gg.vd.pd.distribution.decaylengthxp}{real}{\texttt{0.05}}{Decay length from the X-point for the \texttt{`densitybased'} option.}
\sw{gg.vd.pd.distribution.vdpdincludexp}{integer (0 or 1)}{\texttt{0}}{Include X-points as attraction points.}
\sw{gg.vd.pd.distribution.points.x}{real array}{empty}{x coordinates of additional attraction points (\texttt{`densitybased'} option).}
\sw{gg.vd.pd.distribution.points.y}{real array}{empty}{y coordinates of additional attraction points.}
\sw{gg.vd.pd.distribution.points.d}{real array}{empty}{Decay lengths of additional attraction points.}
\sw{gg.vd.pd.distribution.points.val}{real array}{empty}{Desired density at the additional attraction points.}
\sw{gg.vd.pd.distribution.densityatvessel}{real}{\texttt{250.0}}{Desired density at the vessel boundary.}
\sw{gg.vd.pd.distribution.densityatxp}{real}{\texttt{250.0}}{Desired density at the X-point.}
\sw{gg.vd.pd.distribution.densityatinf}{real}{\texttt{10.0}}{Desired density far from the vessel boundary.}

\subsubsection*{Radial vertex distribution}
\sw{gg.vd.rd.type}{string}{\texttt{`uniform'}}{Same options as the poloidal type (\texttt{`uniform'} or \texttt{`densitybased'}).}
\sw{gg.vd.rd.distribution.doxp}{integer (0 or 1)}{\texttt{0}}{Include X-points (for refinement in the core region).}
\sw{gg.vd.rd.distribution.dosp}{integer (0 or 1)}{\texttt{1}}{Include strike points (for refinement near targets).}
\sw{gg.vd.rd.distribution.fieldwidth}{real}{\texttt{4e-3}}{Desired field width for the uniform distribution.}
\sw{gg.vd.rd.distribution.facelength}{real}{\texttt{0.01}}{Desired face length for the uniform distribution.}
\sw{gg.vd.rd.distribution.decaylengthxp}{real}{\texttt{0.005}}{Decay length parameter at the X-point.}
\sw{gg.vd.rd.distribution.densityatxp}{real}{\texttt{2500.0}}{Desired density at the separatrix / X-point.}
\sw{gg.vd.rd.distribution.decaylengthsp}{real}{\texttt{0.1}}{Decay length parameter at the strike points.}
\sw{gg.vd.rd.distribution.densityatsp}{real}{\texttt{1000.0}}{Desired density at the strike points.}
\sw{gg.vd.rd.distribution.densityatinf}{real}{\texttt{250.0}}{Density far from the separatrix.}
\sw{gg.vd.rd.distribution.points.x}{real array}{empty}{x coordinates of additional attraction points.}
\sw{gg.vd.rd.distribution.points.y}{real array}{empty}{y coordinates of additional attraction points.}
\sw{gg.vd.rd.distribution.points.d}{real array}{empty}{Decay lengths of additional attraction points.}
\sw{gg.vd.rd.distribution.points.val}{real array}{empty}{Desired density at the additional attraction points.}

\subsubsection*{Flux tube extension (cut cells)}
\sw{gg.adap.extendtptubes}{integer (0 or 1)}{\texttt{1}}{Extend tubes at the tangency point side to include (part of) the vessel from both sides of type 1 tangency points.}
\sw{gg.adap.extendvesseltubes}{integer (0 or 1)}{\texttt{0}}{Extend any tube at the vessel edges that complies with the marking criterion.}
\sw{gg.adap.evt.maxvessellength}{real}{\texttt{0.2}}{Maximum L2-based vessel segment length before the tube is extended.}
\sw{gg.adap.evt.noBL}{integer (0 or 1)}{\texttt{1}}{If true, do not apply any boundary layer at extended tubes (typically not desired). Only has effect if boundary layers are applied.}

\subsubsection*{Flux surface removal}
\sw{gg.vd.removefluxsurfaces}{integer (0 or 1)}{\texttt{1}}{(Deprecated) Remove flux surfaces.}
\sw{gg.vd.rfs.mark.psitol}{real}{\texttt{1e-4}}{(Deprecated) Absolute tolerance for the difference in flux value between flux surfaces (if below, one is removed).}
\sw{gg.vd.rfs.mark.psirattol}{real}{\texttt{1e-1}}{(Deprecated) Maximal ratio between flux surface differences (and inverse); if above/below, one is removed.}

\subsubsection*{Boundary triangle removal}
\sw{gg.vd.removenarrowboundarytriangles}{integer (0 or 1)}{\texttt{1}}{Remove narrow boundary triangles.}
\sw{gg.vd.rnbt.mark.criterion}{string}{\texttt{`angle'}}{Criterion for removal; typically \texttt{`angle'} (if too small, remove the triangle).}
\sw{gg.vd.rnbt.mark.minangle}{real}{\texttt{15$\cdot\pi/180$} (rad)}{Minimal angle (too small: removed). The source comment states the input is in degrees, but the default is set in radians.}

\subsubsection*{Small face removal}
\sw{gg.vd.removefaces}{integer (0 or 1)}{\texttt{0}}{Remove small faces.}
\sw{gg.vd.rf.criterion}{string}{\texttt{`facelength\_radial\_bnd'}}{Criterion to remove faces. Typically \texttt{`facelength\_radial\_bnd'} to remove too narrow radial faces at the boundary.}
\sw{gg.vd.rf.minlength}{real}{\texttt{2e-3} [m]}{Minimal length of these faces.}

%======================================================================
\subsection*{Boundary layer (BL) refinement}
%======================================================================
\sw{gg.ref.meth}{string}{\texttt{`no'}}{Refinement method: \texttt{`no'} (no additional refinement) or \texttt{`lengthbased'} (based on min and max length distributions).}
\sw{gg.ref.LB.lengthtype}{string}{\texttt{`euler'}}{Which type of length to consider: \texttt{`euler'} is the classical Eulerian length (L2 norm); \texttt{`radial'} is projected in radial direction (absolute value taken).}
\sw{gg.ref.readexistingrefdata}{integer (0 or 1)}{\texttt{0}}{Read existing refinement data from \texttt{gg.ref.refdatafile}.}
\sw{gg.ref.refdatafile}{string}{\texttt{`'./output/refdataTM.dat'}}{Refinement data file.}
\sw{gg.ref.BL.dotarget}{integer (0 or 1)}{\texttt{0}}{Do BL refinement at targets.}
\sw{gg.ref.BL.dovessel}{integer (0 or 1)}{\texttt{0}}{Do BL refinement at far vessel boundaries.}
\sw{gg.ref.BL.dostructure}{integer (0 or 1)}{\texttt{0}}{Do BL refinement on user-specified structures.}
\sw{gg.ref.BL.nctarget}{integer}{\texttt{0}}{Number of desired boundary layer cells at the target.}
\sw{gg.ref.BL.ncvessel}{integer}{\texttt{0}}{Number of desired boundary layer cells at the vessel.}
\sw{gg.ref.BL.ncstructure}{integer}{\texttt{0}}{Number of desired boundary layer cells at the structures.}
\sw{gg.ref.BL.dltarget}{real array}{empty (fill value \texttt{0.001} [m])}{Desired lengths of the BL cells at the target.}
\sw{gg.ref.BL.dlvessel}{real array}{empty (fill value \texttt{0.01} [m])}{Desired lengths of the BL cells at the vessel.}
\sw{gg.ref.BL.dlstructure}{real array}{empty (fill value \texttt{0.001} [m])}{Desired lengths of the BL cells at the structures.}
\sw{gg.ref.BL.dllengthbased}{integer (0 or 1)}{not set by defaults}{Desired length is specified in classic Euler length or not (if not, the length type is taken).}
\sw{gg.ref.BL.structureID}{integer array}{empty}{Structure-based BL: the structure labels (as specified in the input structure file) at which refinement should be done.}

\subsection*{Length-based refinement (based on exponential decay functions defined in points)}
%======================================================================
\subsubsection*{Poloidal Refinement}
%======================================================================
\sw{gg.ref.LB.doxp}{integer (0 or 1)}{\texttt{1}}{Refine near X-points (X-points are added).}
\sw{gg.ref.LB.dovessel}{integer (0 or 1)}{\texttt{0}}{Refine near (specified) vessel vertices.}
\sw{gg.ref.LB.Lmininf}{real}{\texttt{0.0}}{Minimal length at infinity [m].}
\sw{gg.ref.LB.Lmaxinf}{real}{\texttt{100}}{Maximal length at infinity [m].}
\sw{gg.ref.LB.Lminxp}{real}{\texttt{0.0}}{Minimal length at the X-point [m].}
\sw{gg.ref.LB.Lmaxxp}{real}{\texttt{100}}{Maximal length at the X-point [m].}
\sw{gg.ref.LB.decaylengthxp}{real}{\texttt{0.1}}{Decay length at the X-point (larger: wider influence) [m].}
\sw{gg.ref.LB.Lminstructure}{real array}{empty}{Minimal length on structure $i$ (similar for vertices) [m].}
\sw{gg.ref.LB.Lmaxstructure}{real array}{empty}{Maximal length on structure $i$ [m].}
\sw{gg.ref.LB.decaylengthstructure}{real array}{empty}{Decay length on structure $i$ [m].}
\sw{gg.ref.LB.structureIDs}{integer array}{empty}{Structure IDs of vessel structures to include.}
\sw{gg.ref.LB.vertIDs}{integer array}{empty}{Vessel vertex IDs of vessel vertices to include, as specified in the \texttt{structure.dat} file (some may not be included if deleted during vessel polygon construction).}
\sw{gg.ref.LB.points.x}{real array}{empty}{x coordinates of points $k$ at which poloidal refinement distribution functions are specified.}
\sw{gg.ref.LB.points.y}{real array}{empty}{y coordinates of the points.}
\sw{gg.ref.LB.points.d}{real array}{empty}{Decay lengths at the points [m].}
\sw{gg.ref.LB.points.Lmin}{real array}{empty}{Desired minimal lengths at the points [m].}
\sw{gg.ref.LB.points.Lmax}{real array}{empty}{Desired maximal lengths at the points [m].}

%======================================================================
\subsection*{Radial Refinement}
%======================================================================
\sw{gg.radref.meth}{string}{\texttt{`no'}}{Radial refinement method (same options as the poloidal one).}
\sw{gg.radref.LB.lengthtype}{string}{\texttt{`euler'}}{Type of length to consider (see poloidal option).}
\sw{gg.radref.LB.dosp}{integer (0 or 1)}{\texttt{1}}{Refine near strike points.}
\sw{gg.radref.LB.Lmininf}{real}{\texttt{0.0}}{Minimal length at infinity [m].}
\sw{gg.radref.LB.Lmaxinf}{real}{\texttt{100}}{Maximal length at infinity [m].}
\sw{gg.radref.LB.Lminsp}{real}{\texttt{0.0}}{Minimal length at the strike point [m].}
\sw{gg.radref.LB.Lmaxsp}{real}{\texttt{100}}{Maximal length at the strike point [m].}
\sw{gg.radref.LB.decaylengthsp}{real}{\texttt{0.1}}{Decay length at the strike point (larger: wider influence) [m].}
\sw{gg.radref.BL.dosp}{integer (0 or 1)}{\texttt{0}}{Do BL refinement at strike points.}
\sw{gg.radref.BL.ncsp}{integer}{\texttt{0}}{Number of desired boundary layer cells at the strike point.}
\sw{gg.radref.BL.dlsp}{real array}{empty}{Desired lengths of these cells [m].}
\sw{gg.radref.BL.dllengthbased}{integer (0 or 1)}{not set by defaults}{Desired length is specified in classic Euler length (1) or not (0). If not, the length type is taken.}

\newpage
%======================================================================
\section{Grid Deformation (GD)}
%======================================================================
\sw{gd.main.runtype}{string}{\texttt{`optimize'}}{Type of run for grid deformation. Currently only \texttt{`optimize'} is available; may be extended in the future.}
\sw{gd.main.gridtype}{string}{\texttt{`plasma'}}{Only available option is plasma edge grid (\texttt{`plasma'}).}
\sw{gd.main.meth}{string}{\texttt{`KKT'}}{Method to solve the optimization problem when using \texttt{`optimize'} as runtype. Only \texttt{`KKT'} is supported, though augmented Lagrangian is also available.}
\sw{gd.main.designoptionsfile}{string}{\texttt{`./GOAToptions.dat'}}{File to read design options.}
\sw{gd.main.gridoptionsfile}{string}{\texttt{`./GOAToptions.dat'}}{File to read grid options.}
\sw{gd.main.magneticfieldoptionsfile}{string}{\texttt{`./GOAToptions.dat'}}{File to read magnetic field options.}
\sw{gd.main.numparamsoptionsfile}{string}{\texttt{`./GOAToptions.dat'}}{File to read numerical parameter options.}
\sw{gd.main.environmentoptionsfile}{string}{\texttt{`./GOAToptions.dat'}}{File to read other environment options.}

\newpage
%======================================================================
\section{Grid Adaptation (GA)}
%======================================================================
\subsection*{General}
\sw{ga.plt}{integer (0 or 1)}{\texttt{0}}{Enable general plotting.}
\sw{ga.plt_qm}{integer (0 or 1)}{\texttt{0}}{Enable plotting of quality metrics.}
\sw{ga.meth}{string}{\texttt{`simple'}}{Used method: \texttt{`simple'} adapts the grid based on grid metrics; \texttt{`aposteriori'} uses plasma state information.}

\subsection*{Operations: small and stacked triangles}
\sw{ga.rem_small_trias}{integer (0 or 1)}{\texttt{0}}{Remove small triangles.}
\sw{ga.cut_off_pol}{real}{\texttt{0.3}}{Threshold for small triangles' poloidal height with respect to the poloidal neighbor.}
\sw{ga.cut_off_surf}{real}{\texttt{0.05}}{Threshold for small triangles on the surface area with respect to the surface area of its neighbors.}
\sw{ga.stacked_trias}{integer (0 or 1)}{\texttt{0}}{Apply transformation to stacked triangles.}
\sw{ga.stacked_trias_checkAR}{integer (0 or 1)}{\texttt{0}}{Select triangles based on their aspect ratio.}
\sw{ga.stacked_trias_maxAR}{real}{\texttt{10}}{Maximal aspect ratio of a triangle in the stacked triangle transformation.}
\sw{ga.merge_stacked_trias}{integer (0 or 1)}{\texttt{0}}{Apply merging of stacked triangles which are too skewed.}
\sw{ga.merge_stacked_trias_angle_threshold}{real}{\texttt{5}}{Threshold angle (in degrees) of a stacked triangle to merge it.}
\sw{ga.merge_trap_into_stacked}{integer (0 or 1)}{\texttt{1}}{Allow merging a trapezoid into a group of stacked triangles.}
\sw{ga.stacked_to_cutcell}{integer (0 or 1)}{\texttt{0}}{Apply transformation from stacked triangles to cutcell.}
\sw{ga.stacked_to_cutcell_nonuniform}{integer (0 or 1)}{\texttt{1}}{Transform to cutcell based on the face length of the aligned faces of the triangles.}

\subsection*{Operations: splitting and merging}
\sw{ga.splitting}{integer array}{\texttt{[0]}}{Apply splitting of cells (array of switches, one per splitting pass).}
\sw{ga.merging}{integer array}{\texttt{[0]}}{Apply merging of cells (array of switches, one per merging pass).}
\sw{ga.pents_to_tria}{integer (0 or 1)}{\texttt{0}}{Transform all remaining pentagons into triangles.}
\sw{ga.h_rad_threshold}{real}{\texttt{0.01}}{Threshold of radial width of a cell during merging; only smaller cells are considered.}
\sw{ga.h_rad_core_threshold}{real}{\texttt{0.04}}{Same as \texttt{h\_rad\_threshold}, but only for cells in the core.}

\subsection*{Boundary layer grid (BLG)}
\sw{ga.BLG}{integer (0 or 1)}{\texttt{0}}{Add a boundary layer at the main targets of a grid.}
\sw{ga.BLG_n_layers}{integer}{\texttt{0}}{Number of layers of the added boundary layer.}
\sw{ga.BLG_rescaling_factor}{real}{\texttt{2}}{Size of the boundary layer with respect to the upstream cells.}
\sw{ga.BLG_smoothing_factors}{real array}{\texttt{[1.2, 1.5]}}{The first number is the threshold ratio of subsequent vertex displacement at which a rescaling needs to be done; the second is the maximal ratio allowed (clip).}

\subsection*{Special operations}
\sw{ga.rem_stickout_trias}{integer (0 or 1)}{\texttt{0}}{Remove triangles that are connected to the grid via only one face.}
\sw{ga.rem_trias_tube}{integer (0 or 1)}{\texttt{0}}{Remove flux tubes with only two triangles.}
\sw{ga.rem_outershell}{integer (0 or 1)}{\texttt{0}}{Remove or merge outer flux tubes with a triangle at both ends.}
\sw{ga.rem_tube_outershell_threshold}{real}{\texttt{2}}{Threshold to select a flux tube to be merged with its neighboring tube (h\_rad of neighboring tube / h\_rad of boundary tube).}
\sw{ga.outershell_handling}{string}{\texttt{`merge'}}{Can be \texttt{`remove'} or \texttt{`merge'}.}
\sw{ga.rem_stickout_quad}{integer (0 or 1)}{\texttt{0}}{Remove quads which are connected to the grid with only one face.}
\sw{ga.split_noalignedquads}{integer (0 or 1)}{\texttt{1}}{Split quads without aligned faces.}

\subsection*{Splitting options}
\sw{ga.no_pents}{integer (0 or 1)}{\texttt{1}}{Do not allow pentagons in the final grid.}
\sw{ga.QTtype}{string}{\texttt{`regular'}}{How to handle radial splitting at inclined boundaries: \texttt{`regular'} just splits a triangle radially; \texttt{`pol-rad'} continues poloidal splitting after the radial split of an inclined boundary face.}
\sw{ga.split_out}{integer (0 or 1)}{\texttt{0}}{Option to not split pentagons.}
\sw{ga.splittype}{integer array}{not set by defaults (scalar field \texttt{splittype} defaults to \texttt{`rad'})}{Type of splitting: 1 = radial splitting (reduce radial width of cells); 2 = poloidal splitting (reduce poloidal length of cells).}
\sw{ga.n_split}{integer array}{\texttt{[20]}}{Number of allowed splitting operations.}
\sw{ga.typeT}{string}{\texttt{`cutcell'}}{Method to split a triangle: \texttt{`stacked'} splits it into two triangles; \texttt{`cutcell'} into a quad and a triangle.}
\sw{ga.rad_type}{integer array}{not set by defaults (scalar field \texttt{rad\_type} defaults to \texttt{1})}{Criterion to apply radial splitting (see \texttt{SelectSplitCell}):
1 \texttt{h\_rad\_psi}: split cells with large psi width;
2 \texttt{h\_rad}: split cells with large radial width;
3 \texttt{pol\_flux}: split cells with large poloidal flux estimation based on the \texttt{pol\_flux\_est} distance function;
4 \texttt{farSOL}: split highly inclined triangles in the far SOL;
5 \texttt{farSOL\_refinements}: split cells with large psi width in the far SOL;
6 \texttt{farSOLrefinement\_targets}: split cells with large psi width and low boundary face inclination in the far SOL;
7 \texttt{no\_aligned\_faces}: split cells without aligned faces;
8 \texttt{shaved\_off\_tubes}: split cells in a tube which is concavely shaved off.}
\sw{ga.pol_type}{integer array}{not set by defaults (scalar field \texttt{pol\_type} defaults to \texttt{1})}{Criterion to apply poloidal splitting (see \texttt{SelectSplitCell}):
1 \texttt{trias\_cvS}: split triangles with largest surface area;
2 \texttt{h\_pol}: split cells with largest poloidal length;
3 \texttt{trias\_farSOL}: split triangles with largest surface area in the far SOL;
4 \texttt{farSOLrefinement\_hpol}: split cells with largest poloidal length in the far SOL;
5 \texttt{farSOLrefinement\_targets}: split cells in the far SOL with large poloidal length and low inclination at the boundary face.}
\sw{ga.dist_function_threshold_split}{real}{\texttt{0.9}}{Threshold on the distance function value for selecting a cell to split. The value should be lower than the threshold.}
\sw{ga.dist_function_threshold_split_wall}{real}{\texttt{0.6}}{Same, for the distance function with respect to the wall.}

\subsection*{Merging options}
\sw{ga.no_hex}{integer (0 or 1)}{\texttt{1}}{Do not allow hexagonal cells in the final grid.}
\sw{ga.merge_crit}{integer array}{\texttt{[4]}}{Criterion to select a face over which two cells are merged (see \texttt{SelectMergingFace}):
1 \texttt{tria\_to\_quad}: merge two triangles into quadrilaterals;
2 \texttt{min\_area}: merge cells smaller than the mean surface area;
3 \texttt{h\_pol}: merge cells with too small poloidal length;
4 \texttt{bias}: merge based on high bias between cells in poloidal direction;
5 \texttt{pol\_flux}: merge cells with low poloidal flux estimation via the similarly named distance function;
6 \texttt{h\_rad}: merge cells with psi width smaller than \texttt{h\_rad\_threshold};
7 \texttt{h\_rad\_core}: merge cells in the core with radial width smaller than \texttt{h\_rad\_core\_threshold};
8 \texttt{bias\_rad\_farSOL}: merge cells with large radial bias in the far SOL;
9 \texttt{bias\_rad}: merge cells with large radial bias;
10 \texttt{skew\_tria}: merge triangles with low inclination and high aspect ratio.}
\sw{ga.merge_h_pol_factor}{real}{\texttt{1}}{Factor to loosen the merge criterion based on poloidal length.}
\sw{ga.n_merge}{integer array}{\texttt{[20]}}{Number of merging operations allowed.}
\sw{ga.merge_bias_limit}{real}{\texttt{5}}{Threshold on bias to merge the cells. A merge is done when the real bias is larger than the threshold.}
\sw{ga.dist_function_threshold_merge}{real}{\texttt{0.6}}{Threshold on the distance function value for selecting a face to merge over. The value should be lower than the threshold.}

\subsection*{Pentagon options}
\sw{ga.no_pents_area_merge}{integer (0 or 1)}{\texttt{0}}{Use an area constraint on where pentagons are allowed during merging.}
\sw{ga.no_pents_area_split}{integer (0 or 1)}{\texttt{0}}{Use an area constraint on where pentagons are allowed during splitting.}
\sw{ga.no_pents_area_type}{string}{\texttt{`dist\_function'}}{Method to define the area where no pentagons are allowed: \texttt{`coordinates'} uses the min/max R/Z options below to define a box; \texttt{`dist\_function'} uses a distance function (its value should be lower than the split or merge threshold).}
\sw{ga.no_pents_area_maxR}{real}{\texttt{2.5}}{Maximum R of the box where no pentagons are allowed.}
\sw{ga.no_pents_area_minR}{real}{\texttt{1}}{Minimum R of the box where no pentagons are allowed.}
\sw{ga.no_pents_area_maxZ}{real}{\texttt{-0.9}}{Maximum Z of the box where no pentagons are allowed.}
\sw{ga.no_pents_area_minZ}{real}{\texttt{-2}}{Minimum Z of the box where no pentagons are allowed.}

\subsection*{Distance function options}
\sw{ga.dist_function}{integer (0 or 1)}{\texttt{1}}{Switch to construct distance functions.}
\sw{ga.base_func}{string}{\texttt{`exp(-dist/d)'}}{Basic function for the construction of the distance function. Can only be \texttt{`exp(-dist/d)'}.}
\sw{ga.d_rescale}{real}{\texttt{0.5}}{Rescaling factor for the characteristic length of the distance function.}
\sw{ga.d_rescale_wall}{real}{\texttt{0.5}}{Same as \texttt{d\_rescale}, but for the wall distance function.}
\sw{ga.dist_type}{string}{\texttt{`target\_single\_null'}}{Type of distance function: \texttt{`target\_single\_null'} uses targets as reference; \texttt{`target\_to\_vessel'} uses the outer boundary of the mesh as reference; \texttt{`pol\_flux\_est'} uses the separatrix as reference.}
\sw{ga.dist_type_wall}{string}{\texttt{`target\_to\_vessel'}}{Type of distance function used for the wall distance function (same options as \texttt{ga.dist\_type}).}
\sw{ga.d_char_type}{string}{\texttt{`min\_Xpoint\_dist'}}{Characteristic length used to construct the distance function: \texttt{`min\_Xpoint\_dist'} (minimal distance between targets and X-point) or \texttt{`max\_Xpoint\_dist'} (maximal distance).}
\sw{ga.plt_dist_func}{integer (0 or 1)}{\texttt{0}}{Plot the distance functions.}

\subsection*{A posteriori adaptations}
\sw{ga.vxvol_style}{integer}{\texttt{2}}{Style to interpolate from cell centers to vertices.}
\sw{ga.apost_interpolation_meth}{string}{\texttt{'barycentric'}}{Interpolation method; can be \texttt{`barycentric'} or \texttt{`finite\_element'}.}
\sw{ga.apost_interpolationC}{integer}{\texttt{0}}{Level of continuity of the interpolant.}
\sw{ga.apost_interpolationM}{integer}{\texttt{1}}{Order of the interpolant.}
\sw{ga.apost_meth}{string}{\texttt{`grad'}}{Method to choose a cell to refine: \texttt{`grad'} (from read state and residual information to a newly adapted grid) or \texttt{`res'} (based on residual information).}
\sw{ga.apost_use_na}{integer (0 or 1)}{\texttt{0}}{Use state field \texttt{na} to decide which cells to split.}
\sw{ga.apost_use_ua}{integer (0 or 1)}{\texttt{0}}{Use state field \texttt{ua} to decide which cells to split.}
\sw{ga.apost_use_te}{integer (0 or 1)}{\texttt{0}}{Use state field \texttt{te} to decide which cells to split.}
\sw{ga.apost_use_ti}{integer (0 or 1)}{\texttt{0}}{Use state field \texttt{ti} to decide which cells to split.}
\sw{ga.apost_use_tn}{integer (0 or 1)}{\texttt{0}}{Use state field \texttt{tn} to decide which cells to split.}
\sw{ga.apost_use_po}{integer (0 or 1)}{\texttt{0}}{Use state field \texttt{po} to decide which cells to split.}
\sw{ga.apost_use_kt}{integer (0 or 1)}{\texttt{0}}{Use state field \texttt{kt} to decide which cells to split.}
\sw{ga.apost_use_zt}{integer (0 or 1)}{\texttt{0}}{Use state field \texttt{zt} to decide which cells to split.}
\sw{ga.apost_use_resco}{integer (0 or 1)}{\texttt{0}}{Use residual \texttt{resco} to decide which cells to split.}
\sw{ga.apost_use_resmo}{integer (0 or 1)}{\texttt{0}}{Use residual \texttt{resmo} to decide which cells to split.}
\sw{ga.apost_use_resmt}{integer (0 or 1)}{\texttt{0}}{Use residual \texttt{resmt} to decide which cells to split.}
\sw{ga.apost_use_reshe}{integer (0 or 1)}{\texttt{0}}{Use residual \texttt{reshe} to decide which cells to split.}
\sw{ga.apost_use_reshi}{integer (0 or 1)}{\texttt{0}}{Use residual \texttt{reshi} to decide which cells to split.}
\sw{ga.apost_use_reshn}{integer (0 or 1)}{\texttt{0}}{Use residual \texttt{reshn} to decide which cells to split.}
\sw{ga.apost_use_respo}{integer (0 or 1)}{\texttt{0}}{Use residual \texttt{respo} to decide which cells to split.}
\sw{ga.apost_use_reszt}{integer (0 or 1)}{\texttt{0}}{Use residual \texttt{reszt} to decide which cells to split.}
\sw{ga.apost_use_reskt}{integer (0 or 1)}{\texttt{0}}{Use residual \texttt{reskt} to decide which cells to split.}

\paragraph{Remark.} \texttt{apost\_lambda\_threshold} (default \texttt{3}; cells with length scale per cell width below this threshold are considered for refinement) is documented and has a default, but is not read from the options file.
\end{document}